\section{Conclusion}\label{sec:conclusion}

We prompted whether one spatio-temporal model can collaboratively forecast next-day regional evening peak demand and a graph level Operation Stress Index from Bangladesh national grid datasets under chronological evaluation. The response, on this dataset and methodology is affirmative. Correlation graph PF-STGT with parallel spatial and temporal encoding and repaired multi-task loss provides the finest joint performances in our benchmark suite: macro demand MAE 88.65~MW and OSI $R^2$ 0.745, with statistically substantial demand enhancement relative to the historical hybrid reference.

Three observations stand out. Correlation adjacency approaches better than geographical priors here. Joint training exchanges a small amount of demand accuracy for considerably more useful stress response compared with demand only training. Attribution analysis points to limitation and calendar coalitions for stress and to calendar, generation channels for Dhaka demand, lag, with only limited consistency across methods.

The study introduces a repeatably achievable Bangladesh case for coupled load and stress forecasting and shows that graph design and loss weighting require empirical tuning than hypotheses solely. A single model can enable paired day-ahead forecasting when multi-task calibration, graph structure, and division level monitoring are designed together- with the typical caveat that our findings is from offline evaluation, not from established field deployment. 
